\documentclass[11pt,a4paper]{article}
\usepackage[margin=2.4cm]{geometry}
\usepackage{booktabs,longtable,array,amsmath,microtype}
\usepackage[table]{xcolor}
\usepackage[colorlinks=true,linkcolor=black,urlcolor=blue!55!black]{hyperref}
\newcommand{\pp}{\,\text{pp}}
\newcommand{\gain}[1]{\textcolor{green!45!black}{#1}}
\newcommand{\loss}[1]{\textcolor{red!60!black}{#1}}

\title{\textbf{A Model-Free Explorer for ARC-AGI-3}\\[2pt]
\large Architecture, Ablations, and What the Metric Actually Rewards}
\author{TGAER --- ARC Prize 2026, ARC-AGI-3 track}
\date{29 September 2026}

\begin{document}
\maketitle

\begin{abstract}
\noindent We report on a model-free exploration agent for ARC-AGI-3, a benchmark of
interactive grid games whose rules are never stated. Over the period covered here the
agent was measured against a 25-game local harness with a $0.0300\pp$ noise floor, and
submitted to the Kaggle leaderboard, which scores \emph{110 games the agent has never
seen}. Our headline finding is negative and specific: of fifteen gated mechanism
changes, one shipped with a measurable local gain, one shipped bit-identical and moved
the public score, and thirteen were falsified. We decompose the score to show that
\emph{all} remaining local headroom is action efficiency on games already solved rather
than new games solved, and that the two games holding $64\%$ of that headroom fail for
opposite reasons requiring opposite fixes. We also show that the local benchmark is a
poor proxy in \emph{both} directions, which reframes what a submission should be
selected for.
\end{abstract}

\section{The metric, and why it drives every decision}

A level scores $\min\!\left(\text{cap},\ (h/a)^2\right)$ where $h$ is the human baseline
action count and $a$ is ours. A game averages its levels \emph{weighted by level index
over all of its levels}, so an uncleared level contributes zero to the numerator while
its weight remains in the denominator. Three consequences set the whole research agenda:

\begin{enumerate}\itemsep2pt
  \item \textbf{Quadratic in efficiency.} Being $10\times$ slower than a human costs
  $100\times$ the score. Level \emph{count} is close to irrelevant by comparison.
  \item \textbf{Later levels are worth more.} In an 8-level game level~1 is $1/36$ of
  the score and levels 2--8 are $35/36$. Exploration early is cheap; exploitation late
  is valuable.
  \item \textbf{Actions on a level never cleared are free.} The numerator is zero either
  way, so truncation costs nothing and monotonicity in the action budget holds.
\end{enumerate}

\section{Evaluation protocol}

Local measurements use 25 public games at 600 actions, five seeds per arm, reporting
Relative Human Action Efficiency (RHAE). The gate requires three things, and a candidate
failing any one is rejected: (i) no game may score less often and no game may lose a
level, (ii) wall clock within $1.25\times$ of baseline with both arms projecting inside
the kernel's deadline, and (iii) an argument for why the change should help
\emph{off-roster}. Baseline is $0.3520\% \pm 0.0280$, per-seed
$[0.3827, 0.3827, 0.3316, 0.3316, 0.3316]$; the noise floor is $0.0300\pp$.

Two protocol lessons were learned the hard way and are recorded because they changed
conclusions. First, \textbf{a candidate standard deviation of exactly zero is a defect
report, not a clean result}: the entire seed variance of the 25-game suite comes from
one game (\texttt{lp85}) reaching level 3 on two seeds, so $\sigma\to0$ means that
outcome was pinned. Second, \textbf{the headline verdict is not sufficient}: one
candidate read ``inside noise'' while \texttt{lp85} lost $0.8$ levels on average,
just under the level-drop threshold. The per-game difference must be read every time.

\section{Where the headroom is}

Each game's \texttt{cap} is the score it would post at \emph{perfect} efficiency on the
levels it currently clears. Averaged over the roster this is the ceiling reachable by
pure speed, with no additional level solved. At 6000 actions over three seeds:

\begin{center}
\begin{tabular}{lrrrr}
\toprule
game & levels & env \% & cap \% & capture \\
\midrule
\texttt{lp85} & 3/8 & 3.6330 & 19.444 & 18.7\% \\
\textbf{\texttt{tu93}} & 3/9 & 0.0455 & 13.333 & \textbf{0.3\%} \\
\texttt{m0r0} & 1/6 & 0.0024 & 4.762 & 0.1\% \\
\texttt{sp80} & 1/6 & 0.0345 & 4.762 & 0.7\% \\
\texttt{g50t} & 1/7 & 3.5714 & 3.571 & \textbf{100.0\%} \\
\texttt{ar25} & 1/8 & 1.8701 & 2.778 & 67.3\% \\
\texttt{s5i5} & 1/8 & 0.0006 & 2.778 & 0.0\% \\
\midrule
\multicolumn{2}{l}{\textbf{total}} & \textbf{0.3663} & \textbf{2.0571} & \textbf{17.8\%} \\
\bottomrule
\end{tabular}
\end{center}

\textbf{18 of 25 games score zero, and their \texttt{cap} is zero too}, so coverage work
returns nothing until level~1 actually falls. Everything reachable is efficiency on games
already won. \texttt{g50t} is already at human pace and \texttt{ar25} at $67\%$, so the
two games supplying most of today's score have nothing left to give.

\section{Diagnosis: stranded versus flooded}

Instrumenting the agent's own state graph shows the two games holding $64\%$ of the
ceiling fail in opposite ways.

\begin{center}
\begin{tabular}{lrr}
\toprule
& \texttt{tu93} L2 & \texttt{lp85} L4 \\
\midrule
steps on the level & 5424 / 6000 & 4750 / 6000 \\
\textbf{no route to a frontier} & \textbf{3581 (66\%)} & \textbf{0 (0\%)} \\
current state is itself a frontier & 474 & 3764 (79\%) \\
states discovered & frozen at \textbf{159} & growing to \textbf{1647} \\
untested pairs remaining & frozen at \textbf{6} & growing to \textbf{15950} \\
\bottomrule
\end{tabular}
\end{center}

\texttt{tu93} is \textbf{stranded}: from step 2500 to 6000 it discovers no new state and
drains no untested pair, with six pairs sitting at three states it cannot reach. The
adjacency is directed and holds only observed transitions, so the agent learns
$a \xrightarrow{X} b$ but never the way back unless it happens to take the inverse at
$b$; inside a loop it never does. \texttt{lp85} is \textbf{flooded}: it never fails to
find a route because nearly every state is itself a frontier, and it generates states
faster than it can drain them. A reachability sweep confirms all seven scoring games are
budget-limited rather than out of reachable states.

\section{Ablation ledger}

Every row was gated at five seeds against the noise floor. Rejections are the majority
and are the substance of the work.

\begin{center}
\begin{longtable}{p{6.4cm}rp{5.4cm}}
\toprule
mechanism & $\Delta$ RHAE & outcome \\
\midrule
\endhead
action budget $6000\to12000$ & \gain{$+0.0001\pp$} & shipped; monotone, so free \\
colour-agnostic role inference & $+0.0000\pp$ & \textbf{shipped bit-identical}; moved public $+0.01$ \\
probe cap $4\to1$ & \gain{$+0.1959\pp$} & shipped; largest local gain recorded \\
\midrule
defer irreversible edges & \loss{$-0.1659\pp$} & rejected, $8.4\sigma$; three games stop scoring \\
drop the field-box crop & \loss{$-0.1661\pp$} & rejected; crop is the denoiser \\
re-validate the avatar latch & --- & rejected; \texttt{tu93} $2\to1$ levels \\
$(\text{colour},\text{size})$ compression & --- & not built; candidates collapse only $1.4\times$ \\
class-generalised inert demotion & $+0.0000\pp$ & rejected; bit-identical on all seeds \\
effect (inertness) prediction & \gain{$+0.0065\pp$} & \textbf{closed by ceiling}: $5\times$ under the floor \\
per-level scoping of the stuck test & \loss{$-0.0012\pp$} & rejected; \texttt{tu93} 5/5 seeds $\to$ 0/5 \\
fresh novelty window on respawn & \loss{$-0.0012\pp$} & rejected; \emph{identical} to the row above \\
verify a routed step's destination & $+0.0000\pp$ & bit-identical; routes never mislead \\
learned inverse edges & \loss{$-0.0014\pp$} & rejected; only \texttt{tu93} moved, and it died \\
carry inert evidence across a level & \loss{$-0.0204\pp$} & rejected; $\sigma=0$, \texttt{lp85} pinned \\
\quad same, capped at 1 & \loss{$-0.0204\pp$} & rejected; \emph{bit-identical} to uncapped \\
\quad same, controls only & $+0.0000\pp$ & no-op; 0 of 25 games moved \\
novelty judged against the game's own ceiling & $-0.0002\pp$ & \textbf{submitted}; passes all three gates \\
\bottomrule
\end{longtable}
\end{center}

\subsection{Three results worth more than the fixes would have been}

\paragraph{Effect prediction is closed by a bound, not an experiment.}
A \emph{perfect oracle} for ``this action changes nothing'' is worth
$+0.0065\pp$ against a $0.0300\pp$ floor --- $5\times$ below measurable. The value of
predicting that an action does nothing is bounded by the cost of finding out, which is
one action. Three games have \emph{zero} inert actions in 600 steps. This does not close
mechanisms that avoid \emph{expensive} mistakes.

\paragraph{\texttt{tu93} is a knife edge, so its headroom is not capturable.}
Three independent routing mechanisms each took it from scoring in 5/5 seeds to 0/5, and
the cleanest of them moved \emph{no other game at all}. Its wins depend on an accident:
a death lowers the completed-level count, which reopens a gate while the novelty window
still holds pre-death evidence, so an escape switch fires on the respawned board. Its
nominal $25.9\%$ of the ceiling cannot be reached by changing how it moves, because
reaching its stranded states costs the levels it already has.

\paragraph{Nothing in the inert map transfers across a level.}
Its informative half is clicks, whose keys embed $(\text{row},\text{col})$ and therefore
name a board; its board-free half is 4--6 controls relearned within a handful of steps.
Hence the membership rule: knowledge belongs in a cross-level tier only if its key is
board-free \emph{and} re-earning it is expensive.

\section{Constant sensitivity}

A real mechanism improves over a \emph{range}; a lucky rollout spikes at one value with
baseline either side. Sweeping the module-level tuned constants:

\begin{center}
\begin{tabular}{lll}
\toprule
constant & response & reading \\
\midrule
\texttt{PROBE\_LIMIT} & $0.3825 \to 0.3078 \to 0.3073 \to 0.1865$ & \textbf{monotone} to a floor at 1, not a spike \\
\texttt{CHURN\_FRACTION} & $0.1861,\,0.3800,\,\mathbf{0.3825},\,0.3314,\,0.3314$ & plateau over $0.4$--$0.5$ \\
\bottomrule
\end{tabular}
\end{center}

The sweep tool labelled \texttt{PROBE\_LIMIT} a spike because only one value beats
baseline, but its own definition of a spike requires baseline \emph{on both sides}. The
response is monotone and 1 sits at the boundary; the code further documents that 0 and 1
are identical on all 25 games, with 1 retained to keep the mechanism alive for an unseen
game. The constant is therefore at a natural floor rather than fitted to noise.

\section{Prior art, filtered by the rules we actually play under}

Internet access is disabled in the scored kernel, which are the competition's own terms.
Every published result above ours that uses a frontier API is therefore a research
result on the public set, not a submittable design.

\begin{center}
\begin{tabular}{llll}
\toprule
system & score & model & in-kernel? \\
\midrule
Tycho & 100.0 RHAE, all 183 levels & GPT-5.6 / Opus 5 & \textbf{no} \\
Executable World Models & 58.12\% RHAE & GPT-5.5 & \textbf{no} \\
\textbf{Duck harness} (winner) & \textbf{1.21\%} & Qwen 3.6 27B FP8 & yes \\
AERA & 0.30 (private, 55 games) & Qwen2.5-\textbf{0.5B} & yes \\
\textbf{ours} & \textbf{0.13--0.14} & none (model-free) & yes \\
\bottomrule
\end{tabular}
\end{center}

The benchmark is close to solved by frontier coding agents; the \emph{competition} is
not, because it forbids them. The sharpest datum is not the winner's $9\times$ margin but
that \textbf{AERA beats us roughly $2\times$ using a 0.5B model}, which retires compute
and model size as explanations and leaves architecture. Both kernel-feasible systems
above us are LLM-as-programmer designs. An ablation of the world-model family reports
that \emph{verification} ranks first among its mechanisms and, critically, succeeds at
\emph{lower model effort} --- the regime a local quantised model must work in --- while
the executable world model is ``not universally beneficial''.

\section{The local benchmark is a poor proxy in both directions}

\begin{center}
\begin{tabular}{llll}
\toprule
submission & change & local RHAE & public \\
\midrule
v76 & probe cap $4\to1$ & $0.1561 \to 0.3520\%$ ($2.25\times$) & 0.13 \\
v77 & 12000 actions + colour-agnostic roles & unchanged, \textbf{bit-identical} & \textbf{0.14} \\
repeat & \emph{byte-identical rebuild of v77} & --- & \textbf{0.14} \\
v78 & novelty against the game's own ceiling & $-0.0002\pp$ & pending \\
\bottomrule
\end{tabular}
\end{center}

A $2.25\times$ local gain moved the public score by zero; a local \emph{no-op} moved it
by $+0.01$. The identical rebuild scoring identically is the first evidence that the
public score is reproducible at $0.01$ resolution, which makes v77's $+0.01$ real. The
caveat is that this rests on a single pair, so equality bounds the variance weakly.

The operative conclusion is that a flat local gate is \emph{not} a reason to withhold a
submission, and a local gain is not a reason to make one. What earns a submission is a
named mechanism by which the change differs off-roster. v77 removed a dependency on
public-set colour indices; v78 removes a threshold fitted inside a $0.11$--$0.18$ gap
measured on five of the 25 games that are never scored.

\section{Conclusions}

\begin{enumerate}\itemsep2pt
  \item All remaining local headroom is \textbf{action efficiency on games already
  solved}. Coverage work returns nothing while 18 games score zero.
  \item The two games holding that headroom need \textbf{opposite} fixes, which explains
  why single mechanisms aimed at ``exploration'' kept failing.
  \item \texttt{tu93}'s share of the ceiling is \textbf{not capturable} by routing
  changes, and it has decided three gates on its own. A candidate should be required to
  move a game other than \texttt{tu93}.
  \item \texttt{lp85} is an \textbf{ordering} problem, not a search problem: 15950
  untested pairs where which candidate is tried first is the whole game.
  \item Ordering under a budget is what a value model or a program-writing agent
  supplies, and is the one architecture we have never actually tested. Our earlier 27B
  failure tested LLM-as-\emph{policy}, the weak variant.
\end{enumerate}

\paragraph{Next steps, in order.} Confirm the kernel's runtime limit, which is not stated
in the official documentation while our code assumes 7.5\,h. Extend oracle labels past
level~0 so the transfer metric (actions-per-level regressed on level index) becomes
decisive. Then resume the REPL branches toward the published winner's design, sized to
the confirmed accelerator.

\end{document}
